% source: https://github.com/pfradin/luadraw/discussions/89#discussioncomment-14457144 (pfradin, block 1)
\documentclass{standalone}
\usepackage{xcolor}
\usepackage[3d]{luadraw}
\begin{document}
\begin{luacode*}
------------------
-- spectrum.lua --
------------------

function get_bary_bbox(facets)
  local update_extreme_vals = function (a, a_min, a_max)
    return math.min(a, a_min), math.max(a, a_max)
  end

  local G, level

  local xmin, xmax = math.huge, -math.huge
  local ymin, ymax = math.huge, -math.huge
  local zmin, zmax = math.huge, -math.huge

  for _ , f in ipairs(facets) do
    G = isobar3d(f)

    xmin, xmax = update_extreme_vals(G.x, xmin, xmax)
    ymin, ymax = update_extreme_vals(G.y, ymin, ymax)
    zmin, zmax = update_extreme_vals(G.z, zmin, zmax)
  end

  return {xmin, xmax, ymin, ymax, zmin, zmax}
end

function graph3d:draw_spectrum(surface, spectrum, dir, backcull) -- ajout de l'option backcull=true/false
  local bary_bbox = get_bary_bbox(surface)

  local sorted_facets = self:Sortfacet(surface, backcull) -- Sortfacet de renverra que les facettes "visibles"  si backcull=true

  local imin, imax

  if dir == "x" then
    imin, imax = 1, 2

  elseif dir == "y" then
    imin, imax = 3, 4

  else
    imin, imax = 5, 6
  end

  local val_min   = bary_bbox[imin]
  local delta_val = bary_bbox[imax] - val_min

  local G, val, level

  for _ , f in ipairs(sorted_facets) do
    G = isobar3d(f)

    if dir == "x" then
      val = G.x

    elseif dir == "y" then
      val = G.y

    elseif dir == "z" then
      val = G.z
    end

    level = math.abs((val - val_min) / delta_val)

    self:Dpolyline3d(f, true, "fill=" .. palette(spectrum, level))
  end
end

function graph3d:draw_Xspectrum(surface, spectrum,backcull)
  self:draw_spectrum(surface, spectrum, "x",backcull)
end

function graph3d:draw_Yspectrum(surface, spectrum,backcull)
  self:draw_spectrum(surface, spectrum, "y",backcull)
end

function graph3d:draw_Zspectrum(surface, spectrum,backcull)
  self:draw_spectrum(surface, spectrum, "z",backcull)
end
\end{luacode*}

\begin{luadraw}{name = projective-plane-colorful}
local spectrum = {Gray, SlateGray, LightSkyBlue, LightPink, Pink, LightSalmon}

local pi, sin, cos, tanh = math.pi, math.sin, math.cos, math.tanh

local projplane = surface(
  function(u, v)
    return M(2*(1 + cos(v))*cos(u), 2*(1 + cos(v))*sin(u), 2*sin(v)*tanh(u-pi))
  end,
  pi, 0, 0, pi, -- seulement 1/4 de la figure, avec la bonne orientation
  {25, 15})

-- on complète par symétrie (attention une symétrie inverse l'orientation)
projplane = concat(projplane, reverse_face_orientation(sym3d(projplane,{Origin,vecJ})))
projplane = concat(projplane, reverse_face_orientation(sym3d(projplane,{Origin,vecK})))

local graphview = graph3d:new{
  size    = {12, 12},
  viewdir = {120, 55}
}

graphview:draw_Zspectrum(projplane, spectrum,true) -- avec backcull=true

graphview:Show()
\end{luadraw}

\end{document} 
